% ECONOMICS_PROBLEM: {"id": "J88-323", "title": "Platform-worker protection and flexibility", "jel_code": "J88", "source_id": "OP-0357"}
% !TeX program = xelatex
\documentclass[11pt,a4paper]{article}
\usepackage[margin=23mm]{geometry}
\usepackage{fontspec}
\setmainfont{Latin Modern Roman}
\usepackage{amsmath,amssymb,mathtools}
\usepackage{xcolor,enumitem,booktabs,longtable,array,xurl,hyperref}
\definecolor{muted}{HTML}{53636C}
\hypersetup{colorlinks=true,urlcolor=blue,linkcolor=blue}
\setlength{\parindent}{0pt}
\setlength{\parskip}{5pt}
\newcommand{\R}{\mathbb R}
\newcommand{\E}{\mathbb E}
\newcommand{\N}{\mathbb N}
\newcommand{\ind}{\mathbf 1}
\newcommand{\Var}{\operatorname{Var}}
\newcommand{\Cov}{\operatorname{Cov}}
\newcommand{\supp}{\operatorname{supp}}
\DeclareMathOperator*{\argmin}{arg\,min}
\DeclareMathOperator*{\argmax}{arg\,max}
\newcommand{\entrymeta}[3]{{\sffamily\small\color{muted}\textbf{\texttt{#1}}\quad|\quad #2\par #3\par}}
\newcommand{\Problemref}[1]{\texttt{#1}}
\newcommand{\JELref}[2]{\texttt{#2} (JEL \texttt{#1})}
\begin{document}
\section*{Platform-worker protection and flexibility}
\textbf{Problem ID:} \texttt{J88-323}\quad \textbf{JEL:} \texttt{J88}\par
% BEGIN ECONOMICS STATEMENT
\label{problem:OP-0357}
{\sffamily\small\color{muted}Original subject group: Labor markets and work\par}
\label{definitions:problem:30007639}
\textbf{Audit classification:} Background; proposed extension. Published metadata and full manuscript checked. It supplies platform-surplus and regulatory counterfactuals; the portable-benefits/classification frontier is not its stated open problem.
\subsection*{Definitions and precise research target}
\textbf{Complete editorial problem.} Fix a platform labor market and a policy menu \(\mathcal P\). Policy \(p\) specifies employment classification, an earnings floor, benefit contributions that follow workers across employers, and enforcement. Workers choose participation and hours; buyers choose postings and hires; platforms may change fees or allocation rules. Let \(U_i(p)\) be worker \(i\)'s expected discounted utility, including income risk, benefits, work flexibility, and application costs. Let \(N(p)\) be the normalized share of baseline workers obtaining paid tasks and \(C(p)\) total implementation cost. For explicit welfare weights \(\omega_i\), define \(S(p)=E[\sum_i\omega_iU_i(p)]\). Determine the Pareto frontier of \((S(p),N(p),-C(p))\), including unsuccessful applicants rather than conditioning solely on workers hired. Compare the frontier across plausible demand elasticities and outside employment options. Empirical identification requires policy or contract variation and a model of buyer, worker, and platform responses. A portable-benefits contribution must be traced through fees, wages, prices, and hiring. This is a proposed design problem motivated by existing platform evidence; the cited study does not establish that the entire regulatory menu remains an explicitly stated open problem, and its particular policy results must be preserved.

\textbf{Definitions and admissible scope.} A policy includes rules for unsuccessful applicants and multiple platforms, not just paid hours. Fix a worker probability measure \(\mu\), utility \(U_i(p)=E[\sum_{t=0}^T\delta^t u_i(c_{it},h_{it},a_{it},b_{it})]\), a currency and price base, and finite expectations. Participation and buyer hiring are endogenous. Define \(N(p)=\mu\{i:\text{at least one paid task by }T\}\), and resource cost \(C(p)\) excluding pure transfers that are counted in household budgets. Policy \(p\) dominates \(p'\) if every component of \((S,N,-C)\) is weakly larger and at least one is strictly larger. The frontier consists of nondominated feasible policies. Identification requires an equilibrium response rule and an explicit treatment of cross-platform interference; for a known finite outcome menu frontier computation is elementary. Let the finite worker population be the baseline pool of eligible participants and unsuccessful applicants. Normalize population weights \(\mu_i\) to sum to one and define \(N(p)=\sum_i\mu_i\Pr(\text{at least one paid task by }T\mid p)\), a share; a count multiplies it by the declared pool size. Welfare weights \(\omega_i\), also normalized, may differ from population weights. Costs and utilities have finite means, and a finite feasible policy menu is the base case. Admissible models supply demand, offer allocation, platform pricing, outside options, participation, and an equilibrium selection rule matching the data. The target is the identified set of frontiers over that family. Unknown outcomes, rather than finite-menu frontier computation, create the substantive identification problem.
\subsection*{Variants and assumption changes}
\begin{enumerate}
\item Earnings-floor scope (source-informed specialization): compare a floor per paid hour with a floor per logged-in hour; include waiting, search, and buyer substitution. The cited model supplies particular minimum-wage counterfactuals, not a general portable-benefits theorem.
\item Portable-benefit extension (editorial): impose contribution rate \(\tau\) on gross task revenue and specify vesting and cross-platform portability. Jointly identify incidence on wages, fees, buyer prices, hires, and outside labor supply.
\end{enumerate}
\subsection*{References and source audit}
\textbf{Support and access.} Published metadata and full manuscript checked. It supplies platform-surplus and regulatory counterfactuals; the portable-benefits/classification frontier is not its stated open problem. \textbf{Locator:} November 2024 manuscript, Sections 5--6, pp. 29--32.
\begin{enumerate}\raggedright
\item\hypertarget{definitions:source:30007639:1}{} Christopher T. Stanton; Catherine Thomas. 2025. \emph{Who Benefits from Online Gig Economy Platforms?}. November 2024 manuscript, section 5, printed pp. 29–31; section 6, printed pp. 31–32.
\url{https://www.hbs.edu/ris/Publication%20Files/StantonThomas_Submission3_c26466ef-fcf2-44c6-a356-c92d042113d2.pdf}.
DOI: \url{https://doi.org/10.1257/aer.20221189}.
\end{enumerate}

\subsection*{Common conventions: Source mathematical conventions}

These conventions apply to each entry unless explicitly replaced.

\textbf{Models and identification.} A primitive model \(m\) specifies all probability laws, feasible choices, information, equilibrium and measurement rules used by its target. The admissible class \(\mathcal M\) is fixed independently of outcomes. If \(P_O(m)\) is its observable probability law and \(\tau(m)\) the target, its identified set at observable law \(P\) is
\[
 \mathcal I_\tau(P)=\{\tau(m):m\in\mathcal M,\ P_O(m)=P\}.
\]
Point identification means that this set is a singleton. An empty set signals incompatibility of data and assumptions. Coordinate bounds need not describe the joint set. Bounds are sharp only if every retained target is attainable or arbitrarily approachable by compatible models. Finite bounds require explicit restrictions on unobserved quantities; unrestricted confounding, hidden wealth, extrapolation or arbitrary equilibrium selection can yield uninformative sets.

\textbf{Causal interventions.} A potential outcome \(Y_i(p)\) is the outcome under a complete policy regime \(p\), including timing, financing, information and market spillovers. The target population and units are fixed before treatment. With interference, use the complete assignment vector or a justified exposure mapping; individual no-interference notation alone cannot identify an economy-wide intervention. A difference of mean potential outcomes does not require the cross-world joint distribution. Conditioning on post-treatment migration, survival, enrollment or employment changes the estimand and needs separate justification.

\textbf{Statistical statements.} An estimator, confidence set, forecast or test specifies a sampling law, model class and loss/coverage criterion. Uniform \((1-\alpha)\)-coverage means \(\inf_{m\in\mathcal M}\Pr_m\{\tau(m)\in C_n\}\geq1-\alpha\), or an explicitly asymptotic version. The whole parameter space trivially has coverage, so informativeness requires a stated width, risk or power objective. A forecasting improvement is relative to a named benchmark, information set and evaluation population.

\textbf{Equilibrium and optimization.} An equilibrium comprises feasible plans, optimality, beliefs where needed, budget constraints and all market/resource clearing. The primitives or a stated benchmark define each correspondence \(\mathcal E(p;m)\). Nonemptiness is assumed or proved, never inferred just from compact policy sets. A selected-equilibrium target uses a declared selection rule; a robust target quantifies over all permitted equilibria. Use suprema or infima unless attainment is established. An \(\varepsilon\)-optimal policy is feasible and within \(\varepsilon\) of the finite optimum value. Report an empty feasible set separately.

\textbf{Welfare and accounting.} Normative utility, interpersonal normalization, social weights, numeraire, discounting and the use of public receipts are primitives. Behavioral utility need not coincide with the welfare criterion. Household welfare includes ownership income and rents once; adding profits again to compensating variation generally double-counts them. A monetary equivalent is defined by an explicit transfer and price convention, with monotonicity and range conditions. Average effects, mechanism contributions, transfers and welfare losses are different targets. Infinite sums require convergence and dynamic feasible plans require terminal or no-Ponzi conditions.

\textbf{Source versions.} A working-paper date, revision date and journal publication year are distinguished. A DOI for a journal version is not automatically the DOI of an earlier manuscript. Locators refer to the stated version and printed pages where specified. A metadata-only check does not verify a claim about a full-text conclusion. Every entry's source subsection narrows the provenance claim accordingly.
% END ECONOMICS STATEMENT
\end{document}
